Zorg dat je in de directory LinuxCursus staat en type:

\begin{lstlisting}[language=bash]
$ touch hello.txt
\end{lstlisting}

Na de Enter lijkt er helemaal niets te gebeuren. Dit is met de meeste Linux commando's het geval. Als het goed gegaan is dan laten ze niets weten, een beetje als \textquote{geen nieuws, is goed nieuws}. Doen we een \texttt{ls} dan zien we dat er een bestand is aangemaakt dat hello.txt heet.

Met \texttt{touch}\index{touch} kunnen we dus bestanden aanmaken\index{bestanden maken!leeg}\index{lege bestanden maken}, dit zijn lege bestanden. Type maar eens:

\begin{lstlisting}[language=bash]
$ cat hello.txt
\end{lstlisting}

dan zal je zien dat er weer niets op je scherm verschijnt. En dat is goed! Het \texttt{cat}\index{cat} commando plaatst de inhoud van een bestand op het scherm en daar we een leeg bestand hebben opgevraagd is wat er op het scherm komt dus niets en omdat dat succesvol is verlopen hoeft \texttt{cat} ook geen foutmelding te laten zien en met de wetenschap dat geen nieuws, goed nieuws is is \texttt{cat} klaar.

